\naslov{Sistemi nelineranih enačb}

Rešujemo sistem
\begin{align*}
  f_1(x_1, \ldots, x_n) &= 0 \\
  \vdots& \\
  f_n(x_1, \ldots, x_n) &= 0
\end{align*}
za $f_i : \R \to \R$ ali $\C \to \C$.
Ekvivalentno $F(x) = 0$ za $F: \R^n \to \R^n$ (ali $\C^n \to \C^n$).

Prvi možni pristop reševanja je navadna iteracija.
Sistem $F(x)=0$ zapišemo v ekvivalenti obliki $x = G(x)$, izberemo $x^{(0)}$ ter
iteriramo.

\begin{izrek}
  Naj bo $G: \R^n \to \R^n$ zvezno odvedljiva na zaprti množici $\Omega
  \subseteq \R^n$.
  Če za $x \in \Omega$ velja
  \begin{itemize}
  \item $G(x) \in \Omega$,
  \item $\rho(JG(x)) \le m < 1$,
  \end{itemize}
  kjer je $JG$ Jacobijeva matrika, $\rho$ pa spektralni radij (po absolutni
  vrednosti največja lastna vrednost), potem ima $G$ na $\Omega$ natanko eno
  negibno točko $\alpha$, in za vsak $x^{(0)} \in \Omega$ zaporedje $x^{(r+1)}
  = G(x^{(r)})$ konvergira k $\alpha$.
\end{izrek}

Zadosten pogoj za konvergenco je že, da je $\norm{JG(\alpha)} < 1$ v neki
matrični normi.
Za kvadratično konvergenco mora biti $JG(\alpha) = 0$ po komponentah.

\vprasanje{Kako poiščeš rešitev sistema nelinearnih enačb z navadno iteracijo?
  Povej izrek.}

Podobno kot v enodimenzionalnem primeru lahko uporabimo razvoj v Taylorjevo
vrsto in zanemarimo višje člene.
Dobimo izraz za popravek
\[
  x^{(r+1)} = x^{(r)} - (JF(x^{(r)}))^{-1} F(x^{(r)}).
\]
V praksi raje uporabimo algoritem~\ref{alg:newtonova-metoda}.

\begin{algorithm}
  \caption{Newtonova metoda}
  \label{alg:newtonova-metoda}
  \begin{algorithmic}
	\State Izberi $x^{(0)}$.
	\For{$r = 0, 1, 2, \ldots$}
	\State Reši sistem $JF(x^{(r)}) \Delta x^{(r)} = -F(x^{(r)})$.
	\State $x^{(r+1)} = x^{(r)} + \Delta x^{(r)}$.
	\EndFor
  \end{algorithmic}
\end{algorithm}

\vprasanje{Razloži Newtonovo metodo za rešitev sistema nelinearnih enačb.}

Ker je računanje Jacobijeve matrike zahtevno, se lahko poslužimo kakšne
kvazi-Newtonove metode.
Pri taki metodi na različne načine aproksimiramo Jacobijevo matriko in zmanjšamo
zahtevnost enega koraka.
S tem običajno pade red konvergence na superlinearno.
Najbolj znana kvazi-Newtonova metoda je \pojem{Broydnova metoda}, kjer približek
Jacobijeve matrike $B_{r+1}$ določimo kot najbližjo matriko, ki zadošča
t.i.~\emph{sekantnemu pogoju}
\[
  B_{r+1} (x^{(r+1)} - x^{(r)}) = F(x^{(r+1)}) - F(x^{(r)}).
\]
Ker je $B_r \Delta x^{(r)} = - F(x^{(r)})$, mora torej veljati
\[
  \Delta B_r \Delta x^{(r)} = F(x^{(r+1)}),
\]
matrika $\Delta B_r$ pa je taka, da je $\norm{\Delta B_r}_2$ minimalna.

\begin{lema}
  Dana sta neničelna vektorja $x, y$.
  Matrika $A$ z minimalno normo, ki preslika $x$ v $y$, je
  \[
	A = \frac{y^T x}{\norm{x}_2^2}.
  \]
\end{lema}

\begin{proof}
  Očitno je $Ax = y$.
  Če za matriko $B$ velja $Bx = y$, je $\norm{y}_2 = \norm{Bx}_2 \le \norm{B}_2
  \norm{x}_2$, torej
  \[
	\norm{B}_2 \ge \frac{\norm{y}_2}{\norm{x}_2}.
  \]
  Po drugi strani se da preveriti, da za matrike ranga 1 velja $\norm{y x^T}_2 =
  \norm{y}_2 \norm{x}_2$.
\end{proof}

\begin{algorithm}
  \caption{Broydnova metoda}
  \label{alg:broydnova-metoda}
  \begin{algorithmic}
	\State Določi $x^{(0)}$ in $B_0$.
	\For{$r=0, 1, \ldots$}
	\State Reši $B_r \Delta x^{(r)} = - F(x^{(r)})$.
	\State Izračunaj
	\[
	  B_{r+1} = B_r + \frac{F(x^{(r+1)}) (\Delta x^{(r)})^T}{\norm{\Delta
		  x^{(r)}}_2^2}.
	\]
	\EndFor
  \end{algorithmic}
\end{algorithm}

\vprasanje{Izpelji Broydnovo metodo.}
